\section{Deques}\label{sec:probdeques}

``Deque operations only'' means \cc{push\_front}, \cc{push\_back}, \cc{pop\_front},
\cc{pop\_back}, \cc{front}, \cc{back}, \cc{empty} and \cc{size} of \cc{deque<int>}.

\problem{50}{Design a circular deque}{Medium}
Implement a deque of capacity $k$ in a circular array. The insert and delete operations
return \cc{true} on success and \cc{false} if the deque is full (insert) or empty
(delete); \cc{getFront} and \cc{getRear} return $-1$ if the deque is empty.
\begin{lstlisting}[style=starter]
class CircularDeque {
public:
    int items[1000];
    int cap, first, count;
    CircularDeque(int k);       // called as  CircularDeque d(3);
    bool insertFront(int x);
    bool insertLast(int x);
    bool deleteFront();
    bool deleteLast();
    int getFront();
    int getRear();
    bool isEmpty();
    bool isFull();
};
\end{lstlisting}
\permitted Array access to the fields.
\constraints $1\le k\le1000$.
\required Every operation $O(1)$; no shifting.
\begin{tests}
k=3: insertLast 1, insertLast 2, insertFront 3, insertFront 4 & true, true, true, false\\
getRear, isFull, deleteLast, insertFront 4, getFront & 2, true, true, true, 4\\
k=1: getFront, insertFront 7, insertLast 8, getRear & -1, true, false, 7\\
\end{tests}

\problem{51}{A linked deque with a sentinel}{Medium}
Implement a deque on a circular doubly linked list with one sentinel node:
\cc{sentinel->next} is the first item and \cc{sentinel->prev} the last; when the deque is
empty both point to the sentinel. Write the two helpers first and implement the
operations with them, without any \cc{if} that tests for \cc{NULL}.
\begin{lstlisting}[style=starter]
class dqnode {
public:
    int item;
    dqnode *prev, *next;
    dqnode(int x);
};
class LinkedDeque {
public:
    dqnode *sentinel;
    int count;
    LinkedDeque();                          // creates the sentinel
    void linkBetween(dqnode *a, dqnode *b, int x);
    void unlink(dqnode *n);
    void push_front(int x);   void push_back(int x);
    void pop_front();         void pop_back();
    int front();              int back();
    bool empty();             int size();
    void clear();                           // deletes every item node
};
\end{lstlisting}
\permitted \cc{new}, \cc{delete}, pointer manipulation.
\required Every operation $O(1)$ (\cc{clear}: $O(n)$); no leaks.
\bonus Explain what goes wrong after \cc{LinkedDeque e = d;}.
\begin{tests}
push\_back 1, push\_back 2, push\_front 0 & 0 1 2\\
back, pop\_back; front, pop\_front & 2, 0\\
front, back & 1, 1\\
pop\_back; empty(); pop\_back & (removed); true; nothing happens\\
\end{tests}

\problem{52}{Valid palindrome}{Easy}
Return 1 if the string \cc{s} reads the same forwards and backwards when only its letters
and digits are kept and upper case is treated as lower case; else 0.
\begin{lstlisting}[style=starter]
int isPalindromeText(char s[]);
\end{lstlisting}
\permitted Character reads of \cc{s}; deque operations. No \cc{ctype.h} (write the tests
yourself).
\constraints $1\le|s|\le100$ letters and digits after filtering (the deque has capacity
100).
\required $O(n)$ time.
\bonus $O(1)$ extra space with two indices instead of a deque.
\begin{tests}
A man, a plan, a canal: Panama & 1\\
race a car & 0\\
{" "} (one space) & 1\\
0P & 0\\
Was it a car or a cat I saw? & 1\\
\end{tests}

\problem{53}{Sliding window maximum}{Medium}
For every window \cc{a[i..i+k-1]} write its maximum into \cc{out}; return the number of
windows, $n-k+1$.
\begin{lstlisting}[style=starter]
int maxSlidingWindow(int a[], int n, int k, int out[]);
\end{lstlisting}
\permitted Array access to \cc{a} and \cc{out}; one deque.
\constraints $1\le k\le n\le10^5$, $k\le100$.
\required $O(n)$ time ($O(nk)$ earns partial credit).
\begin{tests}
{}[1,3,-1,-3,5,3,6,7], k=3 & [3,3,5,5,6,7]\\
{}[1], k=1 & [1]\\
{}[9,1,8,2,7,3], k=2 & [9,8,8,7,7]\\
{}[4,2,12,3,8,8,1,5], k=3 & [12,12,12,8,8,8]\\
\end{tests}

\problem{54}{Longest subarray with absolute difference at most a limit}{Medium}
Return the length of the longest non-empty contiguous subarray in which the difference
between the largest and the smallest item is at most \cc{limit}.
\begin{lstlisting}[style=starter]
int longestSubarray(int a[], int n, int limit);
\end{lstlisting}
\permitted Array reads; two deques.
\constraints $1\le n\le10^5$, $0\le\mathit{limit}$; no valid window is longer than 100.
\required $O(n)$ time.
\begin{tests}
{}[8,2,4,7], limit 4 & 2\\
{}[10,1,2,4,7,2], limit 5 & 4 ([2,4,7,2])\\
{}[4,2,2,2,4,4,2,2], limit 0 & 3\\
\end{tests}

\problem{55}{Front-middle-back queue}{Medium}
Design a queue that supports \cc{push} and \cc{pop} at the front, in the middle and at
the back. When there are two middle positions, the \emph{front-most} one is used:
\cc{pushMiddle} inserts so that the new item becomes item $\lfloor n/2\rfloor$ (0-based,
$n$ the size before), and \cc{popMiddle} removes item $\lfloor(n-1)/2\rfloor$. The pops
return the removed item, or $-1$ if the queue is empty.
\begin{lstlisting}[style=starter]
class FrontMiddleBackQueue {
public:
    deque<int> left;
    deque<int> right;
    void pushFront(int x);
    void pushMiddle(int x);
    void pushBack(int x);
    int popFront();
    int popMiddle();
    int popBack();
};
\end{lstlisting}
\permitted Deque operations on the two member deques.
\required Every operation $O(1)$.
\begin{tests}
pushFront 1, pushBack 2, pushMiddle 3, pushMiddle 4 & queue [1,4,3,2]\\
popFront, popMiddle, popMiddle, popBack, popFront & 1, 3, 4, 2, -1\\
\end{tests}

\problem{56}{ArrayDeque: positional insertion and removal}{Medium}
In the ArrayDeque of the course notes, item $i$ ($0\le i<n$) lives at
\cc{a[(j+i) \% 100]}. Implement \cc{add(i, x)}, which inserts $x$ so that it becomes
item $i$, and \cc{remove(i)}, which removes item $i$ and returns it. Each must move the
\emph{shorter} side of the deque: the items before position $i$ (moving \cc{j}) or the
items after it.
\begin{lstlisting}[style=starter]
class ArrayDeque {
public:
    int a[100];
    int j;
    int n;
    ArrayDeque();                // j = 0, n = 0
    int size();
    int get(int i);
    bool add(int i, int x);      // false on failure
    int remove(int i);           // returns the removed item, -1 on failure
};
\end{lstlisting}
\permitted Array access to the fields.
\required $O(1+\min(i,n-i))$ time per operation.
\begin{tests}
add(0,1), add(1,2), add(2,3), add(0,0) & [0,1,2,3]\\
add(2,9) & [0,1,9,2,3]\\
remove(1); get(2) & returns 1, [0,9,2,3]; 2\\
add(5,7); remove(3) & false; returns 3, [0,9,2]\\
\end{tests}

\problem{57}{Jump game VI}{Medium}
You start on index 0 of the array \cc{a} and must reach index $n-1$. From index $i$ you
may jump to any index $i+1,\dots,i+k$ (inside the array). Your score is the sum of
\cc{a[j]} over all indices $j$ you land on, including the first and the last. Return the
maximum score.
\begin{lstlisting}[style=starter]
int maxResult(int a[], int n, int k);
\end{lstlisting}
\permitted Array access; one deque.
\constraints $1\le n\le10^4$, $1\le k\le98$, $|a[i]|\le10^4$.
\required $O(n)$ time ($O(nk)$ earns partial credit).
\begin{tests}
{}[1,-1,-2,4,-7,3], k=2 & 7 (path 1, -1, 4, 3)\\
{}[10,-5,-2,4,0,3], k=3 & 17\\
{}[1,-5,-20,4,-1,3,-6,-3], k=2 & 0\\
\end{tests}
